\chapter*{Resumen}
\addcontentsline{toc}{chapter}{Resumen}

Este Trabajo Fin de Máster estudia dos preguntas que conviene separar. La primera es si un sistema
de aprendizaje automático formado por cinco agentes especializados puede ordenar acciones de mejor
a peor y conservar esa capacidad fuera de muestra. La segunda es cuánto depende el resultado
económico de las reglas que convierten esa ordenación en una cartera. El sistema trabaja con un
universo histórico del S\&P 500, fundamentales fechados por su publicación y cohortes cerradas. Tres
Model Studies encadenados seleccionan el modelo exclusivamente por Rank-IC; después, con la señal
congelada, un Portfolio Study compara \NumCarteras{} configuraciones mediante Information Ratio.

En 2015--2024, el ganador alcanza un Rank-IC medio de \RankICSeleccion{} sobre
\CohortesSeleccion{} cohortes, con \FraccionICPositiva{} positivas, $t$ de Newey--West
\TNeweyWest{} y $p=\PValorPermutacion{}$ en \NumPermutaciones{} permutaciones. La era reservada
2025--2026 mantiene Rank-IC positivo (\RankICReservado{}), pero solo contiene
\CohortesReservadas{} cohortes. Además, el agente \texttt{risk} supera ligeramente al meta-agente,
por lo que la evidencia favorece la capacidad de ordenación, pero no demuestra que la arquitectura
multiagente sea necesaria.

La construcción de cartera altera materialmente el resultado. En selección, el Information Ratio
pasa de \IRModelo{} a \IRCartera{} y la rotación baja de \TurnoverModelo{} a
\TurnoverCartera{} veces al año. En la era reservada, la misma señal pasa de un exceso de
\ExcesoModeloReservado{} con la cartera original a \ExcesoCarteraReservado{} con la optimizada.
Ese cambio de signo es alentador, pero no confirma generalización económica: la cartera es la mejor
de \NumCarteras{} alternativas y la confirmación cubre solo \AnosCarteraReservada{} años. La
contribución principal es, por tanto, metodológica: una señal predictiva y su implementación
económica deben evaluarse por separado y toda conclusión práctica debe declarar la cartera con la
que se midió.

\noindent\textbf{Palabras clave:} aprendizaje automático, selección de acciones, Rank-IC,
\textit{point-in-time}, construcción de cartera, validación fuera de muestra.

\clearpage
\chapter*{Abstract}
\addcontentsline{toc}{chapter}{Abstract}

This Master's Thesis addresses two questions that should not be conflated. First, can a machine
learning system built from five specialised agents rank equities and preserve that ability out of
sample? Second, how strongly does the economic outcome depend on the rules that translate the
ranking into a portfolio? The system uses a historical S\&P 500 universe, filing-dated fundamentals,
and closed cohorts. Three chained Model Studies select the predictive model exclusively through
Rank-IC. Once the signal is frozen, a Portfolio Study compares \NumCarteras{} portfolio
configurations using the Information Ratio.

Over 2015--2024, the selected model obtains a mean Rank-IC of \RankICSeleccion{} across
\CohortesSeleccion{} cohorts, with a Newey--West statistic of \TNeweyWest{} and a permutation
$p$-value of \PValorPermutacion{}. The reserved 2025--2026 period remains positive
(\RankICReservado{}), but contains only \CohortesReservadas{} closed cohorts. The standalone
\texttt{risk} specialist also ranks slightly better than the meta-agent, so the evidence supports
predictive ranking ability without establishing that the multi-agent architecture is necessary.

Portfolio construction materially changes the observed outcome. In the selection window, the
Information Ratio rises from \IRModelo{} to \IRCartera{} while annual turnover falls from
\TurnoverModelo{} to \TurnoverCartera{}. In the reserved period, the same frozen signal changes
from \ExcesoModeloReservado{} excess return under the original portfolio to
\ExcesoCarteraReservado{} under the selected portfolio. The sign reversal is encouraging but does
not establish economic generalisation: the portfolio is the best of \NumCarteras{} alternatives and
the reserved evaluation spans only \AnosCarteraReservada{} years. The main contribution is therefore
methodological: predictive ranking and economic implementation must be evaluated separately, and
any practical claim must state the portfolio through which the signal was measured.

\noindent\textbf{Keywords:} machine learning, equity selection, Rank-IC, point-in-time data,
portfolio construction, out-of-sample validation.
